\documentclass[11pt]{article}
\usepackage[margin=1in]{geometry}
\usepackage{amsmath,amssymb,amsthm,hyperref}
\newtheorem{theorem}{Theorem}
\newtheorem{proposition}[theorem]{Proposition}
\title{Exact two-head sign coalescence and its three-head limit}
\author{}
\date{November 24, 2026}
\begin{document}
\maketitle

\paragraph{Status and scope.}
Theorem~\ref{thm:two} is a complete written statement in the unchanged
affine-ratio model. It has an exact finite checker, but no new Lean
declaration in this note. Proposition~\ref{prop:three} blocks only the
extension that preserves the cleared score values along one prescribed
denominator path. It does not obstruct a path that preserves signs only,
nor either proposed asymptotic bound for worst-case head complexity.

\begin{theorem}[Two-head sign coalescence]\label{thm:two}
Let $n\geq1$, let $A_1,A_2,B_1,B_2$ be affine functions, and suppose
$B_1,B_2$ are positive on the Boolean $n$-cube. Suppose the slopes of
$B_1$ are all strictly positive or all strictly negative. For $0<t\leq1$
put
\[
 C_1=B_1,\qquad C_2(t)=(1-t)B_1+tB_2,
 \qquad N_1(t)=A_1/t,\qquad N_2(t)=A_1+A_2-A_1/t.
\]
Then $C_1,C_2(t)$ are positive on the whole cube, and
\[
 \frac{N_1(t)}{C_1}+\frac{N_2(t)}{C_2(t)}
 =\frac{B_2}{C_2(t)}
   \left(\frac{A_1}{B_1}+\frac{A_2}{B_2}\right).
 \tag{1}
\]
In particular all signs, including zero values, agree. For every
sufficiently small $t>0$, all slopes of $C_2(t)$ have the strict
orientation of $B_1$. Both new atoms are therefore admissible in the
original model, with unrestricted affine numerators, and their
denominator coefficient vectors approach the same coefficient vector.
The hypothesis puts no slope restriction on $B_2$ beyond its original
admissibility, so opposite original orientations are permitted.
\end{theorem}

\begin{proof}
Convexity gives positivity of $C_2(t)$ on the whole cube. The exact
cleared identity is
\[
 \begin{aligned}
 N_1(t)C_2(t)+N_2(t)B_1
 &=\frac{A_1}{t}\big((1-t)B_1+tB_2\big)
   +\left(A_1+A_2-\frac{A_1}{t}\right)B_1\\
 &=A_1B_2+A_2B_1.
 \end{aligned}
\]
Division by $B_1C_2(t)$ proves (1). Its multiplying factor is strictly
positive. If $b_{hi}$ is the slope of $B_h$ in coordinate $i$, then the
new slope is $b_{1i}+t(b_{2i}-b_{1i})$. Finitely many nonzero $b_{1i}$
retain their strict common sign for all sufficiently small positive
$t$. The affine-ratio normal form of the original model admits every
affine numerator when denominator slopes have a strict common sign.
\end{proof}

\paragraph{Bias and constant-denominator cases.}
A free score bias $c$ is absorbed by replacing $A_1$ with $A_1+cB_1$
before the construction. If only one original denominator is
nonconstant, relabel it as $B_1$. If both denominators are constant,
the entire biased score is affine, say $L$; its signs are represented
by the single admissible atom
$L/(1+\varepsilon\sum_i x_i)$ for any $\varepsilon>0$.
Thus every original representation with at most two heads has, at no
increase in head count, a sign-equivalent representation with strict
common denominator orientation and arbitrarily close denominator
coefficient vectors. This is a sign statement; (1) does not preserve
the numerical score itself.

\begin{proposition}[A finite-cube obstruction to fixed-cleared-score
coalescence with three heads]\label{prop:three}
On the three-dimensional Boolean cube write
\[
 B=1+x+y+z,\quad B_2=B+x,\quad B_3=B+y,
 \quad C_2=B+tx,\quad C_3=B+ty.
\]
All these denominators have strictly positive slopes and are positive
on the whole cube whenever $t>0$. For every $t>0$ with $t\neq1$,
there are no affine $L_1,L_2,L_3$ such that
\[
 L_1C_2C_3+L_2BC_3+L_3BC_2=(B+x)(B+y)
 \quad\hbox{on all eight Boolean inputs}.                 \tag{2}
\]
The right side is the cleared score of the original three-head
representation $1/B+0/B_2+0/B_3$.
\end{proposition}

\begin{proof}
Order the cube lexicographically as $000,001,010,011,100,101,110,111$.
Define the linear functional on cube functions by its weight vector
\[
 \lambda_t=\left(
 -(t+1)(t+4),\frac{(t+2)(t+4)}4,
 \frac{t+4}2,-\frac{t+4}3,
 \frac{t+4}2,-\frac{t+4}3,
 -\frac{t+4}{t+3},1\right).
\]
Direct substitution gives, for each $v\in\{1,x,y,z\}$,
\[
 \langle\lambda_t,vC_2C_3\rangle
 =\langle\lambda_t,vBC_3\rangle
 =\langle\lambda_t,vBC_2\rangle=0.                     \tag{3}
\]
These twelve rational identities also hold symbolically in $t$, as
checked independently by the supplied script. A further substitution
gives
\[
 \langle\lambda_t,(B+x)(B+y)\rangle
 =-\frac{(t-1)^2}{t+3}.                               \tag{4}
\]
The right side is nonzero under the stated hypotheses. Affineness and
(3) show that the functional annihilates the left side of (2), whereas
(4) shows that it does not annihilate the right side. This contradiction
proves the claim entirely on the finite cube, allowing all Boolean
identities such as $x^2=x$.
\end{proof}

\paragraph{What remains open.}
The counterexample is deliberately not a sign obstruction: the score
$1/B$ itself retains its positive signs along this denominator path by
keeping its first numerator equal to $1$ and the others zero. More
generally, the same calculation for the nonconstant original score
$(2x-1)/B$ still rules out preserving its original cleared values, but
does not prevent its trivial sign-preserving realization. A many-head
sign-coalescence theorem would have to allow the cleared score to vary,
or use a different denominator path. Neither such a theorem nor a
counterexample to it is proved here.

\paragraph{Verification.}
\texttt{two-head-sign-coalescence-check.py} checks (1) with exact rational
arithmetic on small cubes, including mixed original orientations and
constant second denominators. It checks all twelve identities (3) and
both versions of (4) symbolically, and checks the resulting rank gap at
$t=1/2$. The written proof supplies the universal quantifiers; finite
sampling alone is not used to assert universality.
\end{document}
